\section{Computing the extremal upper density}
There is an integer precedent for computing extremal upper densities by
optimizing over a finite set of primes and then enlarging it.
See N. McNew, \emph{On sets of integers which contain no three terms in
geometric progression}, Math. Comp. 84 (2015), 2893--2910,
\href{https://arxiv.org/abs/1310.2277}{Theorems 4.2 and 5.1}.
The earlier incoming \texttt{triangle.tex} uses one and two linear primes
to bound the supremum over polynomials. The following adaptation to all
irreducibles of degree at most $T$ gives a bound on the approximation error
in terms of the degree cutoff.
Fix $q$. Write $\beta$ for the supremum in Conjecture 6.3.
For an integer $T\ge1$, let $S$ comprise all monic irreducibles of degree at
most $T$, and let $J_S(t)$ be the largest size of a progression-free subset
of the monic $S$-smooth polynomials of degree at most $t$. Put
\[
 U_T=(1-q^{-1})\prod_{P\in S}(1-q^{-\deg P})
              \sum_{t\ge0}J_S(t)q^{-t},\qquad
 v_T=\frac{1-q^{-1}}{\prod_{P\in S}(1-q^{-2\deg P})}.
\]
Then
\[
 v_T U_T\le\beta\le U_T,\qquad
 0\le U_T-\beta\le 1-v_T\le \frac{q^{-T}}{q-1}.
\]
Thus $\beta=\lim_{T\to\infty}U_T$ is a computable real. The formula
also holds with rational-function ratios if that convention is used for
both $\beta$ and every $J_S(t)$.

For the upper bound, fix the scalar and the monic rough part $g$.
At cutoff $N$, at most $J_S(N-\deg g)$ smooth factors can occur in that
class. For the lower bound, restrict the allowed $g$ to be squarefree
and use a maximizing subset independently in each class, using the same
choice for all scalars at fixed $g$. In any allowed progression the valuations
outside $S$ are three equally spaced members of $\{0,1\}$, hence constant.
The three monic parts therefore lie in one class, where their quotients
would give the prohibited progression between smooth factors. This also
handles signed valuation steps in the rational-ratio version, including
nonconstant ratios of degree zero. To glue the finite constructions,
use the degree blocks from Section~\ref{sec:density}. A rational-ratio
progression with nonconstant degrees can be ordered in increasing degree
by reversing its ratio if necessary. Equal-degree triples are already
confined to one block.

The proportion of monic degree-$n$ polynomials coprime to all primes in
$S$ tends to $c=\prod_{P\in S}(1-q^{-\deg P})$, by finite
inclusion--exclusion. The proportion also squarefree tends to $cv_T$ by
the generating function in Section~\ref{sec:density}. In the normalized
count, the classes with $\deg g=N-t$ therefore contribute at most
$(1-q^{-1})c\,J_S(t)q^{-t}$ in the limit, with the corresponding
lower construction attaining the same expression times $v_T$. Dominated
convergence applies since $J_S(t)\le(t+1)^{|S|}$ by counting the possible
exponent vectors. Also $U_T\le1$, by applying the identical count with
all smooth factors allowed. Finally
\[
 1-v_T=1-\prod_{P\notin S}(1-q^{-2\deg P})
 \le \sum_{d>T}q^d q^{-2d}=\frac{q^{-T}}{q-1}.
\]
For each fixed $T$, finite enumeration of admissible sets computes
$J_S(t)$ exactly. Truncating its series at $t=L$ costs at most the
explicit tail $\sum_{t>L}(t+1)^{|S|}q^{-t}$, which tends effectively to zero.
This characterizes the supremum to arbitrary prescribed error. Obtaining
a practically sharper numerical interval for small $q$ requires solving
large finite extremal problems.

The supremum is attained as an upper density. For stage $j$ take $T=j$,
and choose $N_j$ sufficiently large that the squarefree construction has
density at least $v_jU_j-1/j$ in its ball and satisfies both block-separation
inequalities of Section~\ref{sec:density}. Keep its top $j+1$ degree layers.
Their union is progression-free and has upper density at least
$\lim_{j\to\infty}(v_jU_j-1/j-q^{-j-1})=\beta$. By definition of
$\beta$ equality follows.
